\setcounter{equation}{0}
\chapter{Introduction}

In particular to this thesis is the study of \quotes{complex fluids} or \quotes{soft materials} which blur the traditional line between solids and liquids. These classes of material, which encompass polymers, emulsions, foams and liquid crystals, amongst many others, have complex microstructure and exhibit more solid-like properties, whilst still displaying the typical properties associated with fluids. Many Non-Newtonian fluids have been shown to demonstrate the formation of shear banded deformations \cite{BerretWormlikeMicelles2004,bonn_yield_2017,mfielding_shear_2009,lockwood_ultradelayed_2025, nicolas_deformation_2018,pollard_yielding_2022, radhakrishnan_shear_2016, radhakrishnan_shear_2018,cochran_slow_2024,martens_spontaneous_2012,fielding_complex_2007,fielding_shear_2014,manneville_recent_2008,OlmstedShearBanding2008}. Shear banding in the context of this thesis can be defined as a regime where the material's strain, or deformation, becomes highly concentrated into a narrow, distinct region during large plastic deformation.


\section{Layout of thesis}

\section{Notation}
\subsection*{Vectors, tensors, and scalars}

\begin{itemize}
    \item Bold symbols denote vectors and tensors. For example, the position vector is written as $\mathbf{r}$ and the stress tensor as $\mathbf{\Sigma}$.

    \item Non-bold symbols denote scalar quantities. For example, the volume fraction is written as $\phi$.

    \item Given a bold symbol, the corresponding non-bold symbol, in the absence of a sub- or superscript specifying a component, denotes its magnitude. For example,
    \[
        r := |\mathbf{r}|.
    \]
\end{itemize}

\subsection*{Relative vectors}

Bold symbols with a double subscript denote relative quantities between nodes $i$ and $j$. For example, the relative position of node $j$ with respect to node $i$ is defined as
\[
    \mathbf{r}_{ij} := \mathbf{r}_j - \mathbf{r}_i.
\]

\subsection*{Index notation}

Throughout this thesis, we adopt the Einstein summation convention, whereby repeated indices are implicitly summed over.

\noindent Partial derivatives with respect to Cartesian coordinates are written as
\[
    \partial_i \equiv \frac{\partial}{\partial x_i}.
\]

\noindent The Kronecker delta is defined as
\begin{equation}
    \delta_{ij} =
    \begin{cases}
        1, & \text{if } i=j,\\
        0, & \text{if } i\neq j,
    \end{cases}
    \label{symbol:delta}
\end{equation}
and acts as the identity tensor in index notation.

\subsection*{Fourier-space notation}

A circumflex (``hat'') denotes the Fourier transform of a quantity. For a function $f(x)$, we write
\[
    \hat{f}(k) = \mathcal{F}(f)(k),
\]
where $\mathcal{F}$ denotes the Fourier transform and $k$ is the wavenumber.

\subsection*{Special functions}

The Dirac delta distribution is defined implicitly through
\begin{equation}
    \int_{-\infty}^{\infty} \delta(x)\, \mathrm{d}x = 1,
    \qquad
    \int_{-\infty}^{\infty} f(x)\,\delta(x-x_0)\,\mathrm{d}x = f(x_0),
    \label{symbol:dirac_delta}
\end{equation}
for any sufficiently smooth function $f(x)$.

\noindent The Heaviside step function is defined as
\begin{equation}
    \Theta(x)= 
    \begin{cases} 
        0, & \text{if } x < 0, \\ 
        1, & \text{if } x \ge 0,
    \end{cases}
    \label{symbol:heaviside_function}
\end{equation}
which is used throughout this thesis to represent discontinuous transitions.